%%% XeLaTeX-article %%%
%# -*- coding: utf-8 -*-
%!TEX encoding = UTF-8 Unicode
%!TEX TS-program = xelatex
%---------------------虽然加了%还是要保留！

\documentclass[17pt]{beamer}
\mode<presentation>
{
\usetheme[width=40pt]{Hannover}
\usecolortheme[]{dove}
\usefonttheme[]{structurebold}
\setbeameroption{hide notes}
}

\usepackage{fontspec}
\usepackage{polyglossia}
\setmainfont{Arial} %设置主字体
\newfontfamily\sanskritfont[Script=Devanagari,Mapping=romantodevanagari,Scale=1.15]{Sanskrit 2003}             %输出天城体
%\newfontfamily\sanskritfont[Mapping=tex-text]{Times New Roman}              %输出转写
\doublehyphendemerits=-10000
\newcommand{\skt}[1]{{\sanskritfont{#1}}} %输出天城体
\newcommand{\skttrans}[1]{{\skt{#1}~#1}}  %输出天城体和转写
%----------------------------------------------------设置梵文输入方法 danda । ॥

\usepackage[UTF8,fontset=windows]{ctex}
\usepackage{amsmath}
%----------------------------------------------------设置中文环境

\usepackage{graphicx}
\usepackage{flafter}
\graphicspath{{pic/}}
\usepackage{booktabs}
\usepackage{nicematrix}
\newenvironment{indentlist}
  {\begin{list}{}{\setlength{\leftmargin}{2em}\setlength{\itemsep}{0.5em}}}
  {\end{list}}
%-----------------------------------------插图表格

\usepackage{hyperref}
\usepackage[dvipsnames]{xcolor}
\usepackage{colortbl}
\definecolor{light-gray}{gray}{0.9}
%------------------------------颜色

\newcommand{\verbroot}[1]{\textcolor{red}{$\sqrt{}$#1}}
\newcommand{\sktroot}[1]{{\verbroot{\skt{#1}}}}
\newcommand{\skttransroot}[1]{{\sktroot{#1}~\textcolor{red}{#1}}}

\newcommand{\nounstem}[1]{\textcolor{red}{#1\nobreakdash-}}
\newcommand{\sktnounstem}[1]{{\textcolor{red}{\skt{#1\nobreakdash-}}}}
\newcommand{\skttransnounstem}[1]{{\sktnounstem{#1}~\nounstem{#1}}}

\newcommand{\verbstem}[1]{\textcolor{blue}{#1\nobreakdash-}}
\newcommand{\sktverbstem}[1]{{\textcolor{blue}{\skt{#1\nobreakdash-}}}}
\newcommand{\skttransverbstem}[1]{{\sktverbstem{#1}~\verbstem{#1}}}

\newcommand{\wordending}[1]{\textcolor{Orange}{\nobreakdash-#1}}
\newcommand{\sktending}[1]{{\textcolor{Orange}{\skt{-#1}}}}
\newcommand{\skttransending}[1]{{\sktending{#1}~\wordending{#1}}}

\newcommand{\fullpada}[1]{\textcolor{OliveGreen}{#1}}
\newcommand{\sktpada}[1]{{\textcolor{OliveGreen}{\skt{#1}}}}
\newcommand{\skttranspada}[1]{{\sktpada{#1}~\fullpada{#1}}}

\newcommand{\pratyaya}[1]{\mbox{\textcolor{Plum}{#1}}}
\newcommand{\sktpratyaya}[1]{{\textcolor{Plum}{\skt{#1}}}}
\newcommand{\skttranspratyaya}[1]{{\sktpratyaya{#1}~\pratyaya{#1}}}

\newcommand{\reconstruction}[1]{\textcolor{gray}{*#1}}
\newcommand{\fullsentence}[1]{\textcolor{MidnightBlue}{#1}}
\newcommand{\sktsentence}[1]{\textcolor{MidnightBlue}{\skt{#1}}}

\newcommand{\veryimportant}[1]{\textcolor{red}{#1}}
\newcommand{\important}[1]{\textcolor{blue}{#1}}
\newcommand{\notsoimportant}[1]{\textcolor{gray}{#1}}
\newcommand{\dicturl}[2]{\href{#1}{\mbox{\texttt{#2}}}}
%-------------------------------------------词根等标颜色

\title{{梵语提高}}
\subtitle{38. 不定过去时}
\author[张雪杉]{文学院~~张雪杉\\zhangxueshan@sdnu.edu.cn}
\date{}

\begin{document}

\begin{frame}
  \titlepage
\end{frame}

\begin{frame}
  \frametitle{本节内容}
  \small
  \tableofcontents
\end{frame}

\section{上节作业}

\begin{frame}{《薄伽梵歌》1.34--37}
  \footnotesize{
    \begin{verse}
      \mbox{\skt{ācāryāḥ pitaraḥ putrāstathaiva ca pitāmahāḥ ।}}\\
      \mbox{\skt{mātulāḥ śvaśurāḥ pautrāḥ śyālāḥ saṃbandhinastathā ।। 1-34 ।।}}\\
      \skt{etānna hantumicchāmi ghnato 'pi madhusūdana ।}\\
      \mbox{\skt{api trailokyarājyasya hetoḥ kiṃ nu mahīkṛte ।। 1-35 ।।}}\\
      \skt{nihatya dhārtarāṣṭrānnaḥ kā prītiḥ syājjanārdana ।}\\
      \mbox{\skt{pāpamevāśrayedasmānhatvaitānātatāyinaḥ ।। 1-36 ।।}}\\
      \skt{tasmānnārhā vayaṃ hantuṃ dhārtarāṣṭrānsvabāndhavān ।}\\
      \mbox{\skt{svajanaṃ hi kathaṃ hatvā sukhinaḥ syāma mādhava ।। 1-37 ।।}}\\
    \end{verse}
  }
\end{frame}

\begin{frame}{《薄伽梵歌》1.38--40}
  \footnotesize{
    \begin{verse}
      \skt{yadyapyete na paśyanti lobhopahata-cetasaḥ ।}\\
      \mbox{\skt{kulakṣayakṛtaṃ doṣaṃ mitradrohe ca pātakam ।। 1-38 ।।}}\\
      \skt{kathaṃ na jñeyamasmābhiḥ pāpādasmānnivartitum ।}\\
      \mbox{\skt{kulakṣayakṛtaṃ doṣaṃ prapaśyadbhirjanārdana ।। 1-39 ।।}}\\
      \skt{kulakṣaye praṇaśyanti kuladharmāḥ sanātanāḥ ।}\\
      \mbox{\skt{dharme naṣṭe kulaṃ kṛtsnamadharmo 'bhibhavatyuta ।। 1-40 ।।}}\\
    \end{verse}
  }
\end{frame}

\section[不定过去]{不定过去时}
\begin{frame}{\insertsection}
  \small
  \tableofcontents[currentsection]
\end{frame}

\subsection{概述}
\begin{frame}{表示过去的时态}
  \small
  \begin{itemize}
    \item 古典梵语三种过去时含义相同：\\

    \item 未完成时（第12章）：\\
    前加元音 \pratyaya{a-} + \important{现在时语干} + 派生语尾
    \item 完成时（第27--28章）：\\
    \mbox{重复音节 + \important{等级完成时语干} + 完成时语尾}
    \item \notsoimportant{迂回完成时（第34章）：\\
    \mbox{现在时语干 + \pratyaya{-ām} + 助动词}}
    \item \veryimportant{不定过去时}：\\
    \mbox{前加元音 \pratyaya{a-} + \important{不定过去时语干} + 派生语尾}
  \end{itemize}
\end{frame}

\begin{frame}{\insertsection}
  \small
  \begin{itemize}
    \item 不定过去时最少见。
    \item 识别重点：\\
    \mbox{前加元音 \pratyaya{a-} + 派生语尾：和未完成时一样}\\
    \veryimportant{非现在时语干} ：和以前见过的哪里不一样
    \item 最常见用法：\\
    \fullpada{mā} + \veryimportant{不含前加元音 \pratyaya{a-} }的不定过去时\\
    表示禁止的命令
  \end{itemize}
\end{frame}

\subsection{七种类型}

\begin{frame}{七种类型例表}
  \footnotesize
  \centering
  \resizebox{\textwidth}{!}{%
  \begin{NiceTabular}{|c|c|c|c|c|c|c|c|}[hvlines, rules/width=0.3pt, rules/color=gray]
    & \important{1a) bhū} & \important{1b) vid} & \important{2) jan}
    & \important{3a) rudh} & \important{3b) budh} & \important{3c) yā} & \important{3d) diś} \\
    Sg.1 & \fullpada{abhūvam} & \fullpada{avidam} & \fullpada{ajījanam}
         & \fullpada{arautsam} & \fullpada{abodhiṣam} & \fullpada{ayāsiṣam} & \fullpada{adikṣam} \\
    Sg.2 & \fullpada{abhūḥ} & \fullpada{avidaḥ} & \fullpada{ajījanaḥ}
         & \fullpada{arautsīḥ} & \fullpada{abodhīḥ} & \fullpada{ayāsīḥ} & \fullpada{adikṣaḥ} \\
    Sg.3 & \fullpada{abhūt} & \fullpada{avidat} & \fullpada{ajījanat}
         & \fullpada{arautsīt} & \fullpada{abodhīt} & \fullpada{ayāsīt} & \fullpada{adikṣat} \\
    Du.1 & \fullpada{abhūva} & \fullpada{avidāva} & \fullpada{ajījanāva}
         & \fullpada{arautsva} & \fullpada{abodhiṣva} & \fullpada{ayāsiṣva} & \fullpada{adikṣāva} \\
    Du.2 & \fullpada{abhūtam} & \fullpada{avidatam} & \fullpada{ajījanatam}
         & \fullpada{arautttam} & \fullpada{abodhiṣṭam} & \fullpada{ayāsiṣṭam} & \fullpada{adikṣatam} \\
    Du.3 & \fullpada{abhūtām} & \fullpada{avidatām} & \fullpada{ajījanatām}
         & \fullpada{arautttām} & \fullpada{abodhiṣṭām} & \fullpada{ayāsiṣṭām} & \fullpada{adikṣatām} \\
    Pl.1 & \fullpada{abhūma} & \fullpada{avidāma} & \fullpada{ajījanāma}
         & \fullpada{arautsma} & \fullpada{abodhiṣma} & \fullpada{ayāsiṣma} & \fullpada{adikṣāma} \\
    Pl.2 & \fullpada{abhūta} & \fullpada{avidata} & \fullpada{ajījanata}
         & \fullpada{arautta} & \fullpada{abodhiṣṭa} & \fullpada{ayāsiṣṭa} & \fullpada{adikṣata} \\
    Pl.3 & \fullpada{abhūvan} & \fullpada{avidan} & \fullpada{ajījanan}
         & \fullpada{arautsuḥ} & \fullpada{abodhiṣuḥ} & \fullpada{ayāsiṣuḥ} & \fullpada{adikṣan} \\
  \end{NiceTabular}%
  }
\end{frame}

\subsection{简单不定过去时}
\begin{frame}{简单不定过去时}
  \begin{itemize}
    \item[1a)] \important{词根不定过去时}：\\
    词根直接加语尾 \\
    \notsoimportant{限 \verbroot{bhū} 、复合元音与 ā 结尾词根}
    \item[1b)] \important{a-不定过去时}：\\词根零级 + \pratyaya{-a-}~~\notsoimportant{（最常见）}
    \item[2)] \important{重复不定过去时}：\\重复语干 + \pratyaya{-a-}~~\notsoimportant{（派生动词）}
  \end{itemize}
\end{frame}



\begin{frame}{1b~~a-不定过去时}
  \small
  \begin{itemize}
    \item 语干构成：\important{词根零级} + \pratyaya{-a-} 
    \item 许多一四六类动词遵从这种形式：\\
    \verbroot{gam} (1) 'go'  Impf. \fullpada{agacchat} Aor. \fullpada{agamat}\\
    \verbroot{śuc} (1) 'grieve'  Impf. \fullpada{aśocat} Aor. \fullpada{aśucat}\\
    \verbroot{vid} (6) 'know'  Impf. \fullpada{avindat} Aor. \fullpada{avidat}\\
  \end{itemize}
\end{frame}

\begin{frame}{2~~重复不定过去时}
  \small
  \begin{itemize}
    \item 语干构成：\important{重复语干} + \pratyaya{-a-} 
    \item 重复规则：\\
    u 重复 u/ū，其他重复 i/ī\\
    重复音节和词根音节长短不同
    \item 主要用于\important{派生动词}（致使、第十类动词）\\
    \verbroot{jan} 'produce' \\ Caus. Impf. \fullpada{ajanayat} Aor. \fullpada{ajījanam}\\
    \verbroot{jīv} 'make alive'  Caus. \fullpada{jīvayati} Aor. \fullpada{ajījivat}\\
    \verbroot{budh} 'wake up' \\ Caus. \fullpada{bodhayati} Aor. \fullpada{abūbudhat}\\
  \end{itemize}
\end{frame}

\subsection{咝音不定过去时}
\begin{frame}{3~~咝音不定过去时}
  \begin{itemize}
    \item[3a)] \important{s-不定过去时}~~
    \item[3b)] \important{iṣ-不定过去时}~~
    \item[3c)] \notsoimportant{siṣ-不定过去时}
    \item[3d)] \notsoimportant{sa-不定过去时（后两种少见）}
  \end{itemize}
\end{frame}

\begin{frame}{3a~~s-不定过去时}
  \small
  \begin{itemize}
    \item 语干构成：词根元音变化 + \pratyaya{-s-} 
    \item \important{主动语态}：词根三合\\
    \important{中间语态}：词根零级\notsoimportant{（i/u结尾二合）}\\
    \verbroot{rudh} P. \fullpada{arautsīt} Ā. \fullpada{aruddha} \\
    \notsoimportant{注意连声，\pratyaya{-s-} 有时还会脱落}
    \item 词尾变化：\\
    P.2.Sg. \pratyaya{-īḥ} P.3.Sg. \pratyaya{-īt} \\
    P.3.Pl. \pratyaya{-uḥ} \notsoimportant{（不带插入元音 \pratyaya{-a-} 的都这样）}
  \end{itemize}
\end{frame}



\subsection{mā 表禁止}

\begin{frame}{\insertsubsection}
  \small
  \begin{itemize}
    \item 不定过去时古典梵语最常见的用法：\\
    \fullpada{mā} + \veryimportant{不含前加元音 \pratyaya{a-} }的不定过去时\\
    表示禁止某事（否定命令）
    \item 例：\\
    \fullsentence{mā bhūḥ} \notsoimportant{"别（这样）存在/做！"}\\
    \verbroot{bhū}，词根不定过去时 2.Sg.    \\
    \fullsentence{mā śucaḥ} \notsoimportant{"不要悲伤！"}\\
    \verbroot{śuc}，a-不定过去时 2.Sg.
    
  \end{itemize}
\end{frame}


\section{本节作业}

\begin{frame}{\insertsection}
  \begin{itemize}
    \item 第35章阅读练习
  \end{itemize}
\end{frame}

\end{document}
